% Part: sets-functions-relations
% Chapter: size-of-sets-complete

\documentclass[../../../include/open-logic-chapter]{subfiles}

\begin{document}

\olchapter{sfr}{siz}{اندازهٔ مجموعه‌ها}

\begin{editorial}
این فصل به برشماری‌ها، شمارایی و ناشمارایی می‌پردازد. چند بخش در دو
روایت آمده‌اند: روایتی مقدماتی‌تر که برشماری‌ها را فهرست‌ها، یا توابع
پوشا از $\PosInt$، در نظر می‌گیرد؛ و روایتی انتزاعی‌تر که برشماری‌ها را
تناظرهای دوسویی با $\Nat$ تعریف می‌کند.
\end{editorial}

\olimport{introduction}

\olimport{enumerability}

\olimport{zig-zag}

\olimport{pairing}

\olimport{pairing-alt}

\olimport{non-enumerability}

\olimport{reduction}

\olimport{equinumerous-sets}

\olimport{comparing-size}

\olimport{schroder-bernstein}

\begin{editorial}
بخش‌های زیر، یعنی \olref[sfr][siz][enm-alt]{sec}،
\olref[sfr][siz][nen-alt]{sec} و \olref[sfr][siz][red-alt]{sec}،
روایت‌های جایگزینی از \olref[sfr][siz][enm]{sec}،
\olref[sfr][siz][nen]{sec} و \olref[sfr][siz][red]{sec} هستند که تیم
باتن برای استفاده در متن خود با عنوان «نظریهٔ مجموعه‌های باز» نوشته
است. این روایت‌ها اندکی پیشرفته‌ترند و تعریفی متفاوت از
شمارش‌پذیری را به کار می‌گیرند که برای بافت نظریهٔ مجموعه‌ها مناسب‌تر
است (یعنی تناظر دوسویی با $\Nat$ یا یک پارهٔ آغازین، نه
فهرست‌پذیر بودن یا برد تابعی پوشا از $\PosInt$ بودن).
\end{editorial}

\olimport{enumerability-alt}
\olimport{non-enumerability-alt}
\olimport{reduction-alt}

\OLEndChapterHook

\end{document}
